\section{Conclusions}\label{sec:conclusions}

The spring-coupled inverted pendulum has been developed into a transparent benchmark for the recovery of missing dynamics and used to extend the full--partial reconstruction mapping with a generative engine. Five conclusions emerge.

First, the conservative dynamics is governed by the single parameter $\mu=kL/(Mg)-1$, which unfolds a subcritical pitchfork at $\mu=0$. Near this critical point the resonance frequency and the tilt of the saddles vanish as $\mu^{1/2}$ but the barrier vanishes as $\mu^{2}$, which quantifies the trade-off between isolation and robustness that governs inverted-pendulum isolators. Five canonical regimes, computed with trajectories that conserve energy to $10^{-10}\,MgL$ and periods that agree with exact quadratures, display a libration inside the upright well, whose period is lengthened by $4.2\%$, in agreement with the $4.1\%$ predicted by the normal form; critical slowing down at the degenerate point; and rotations near and far from the separatrix.

Second, under damping and forcing the separatrix splits and seeds chaos: no chaotic state was found below the Melnikov threshold curve $A_c(\Omega)$, which gives $(A/\gamma)_c=2.049$ at the reference drive frequency, and the dynamics reaches, through a period-doubling cascade with ratios converging to Feigenbaum's constant, a strange attractor with $\lambda_1=0.0956\,(g/L)^{1/2}$ and $D_{\mathrm{KY}}=2.277$, whose dimension is consistent with the embedding dimension $d_E=5$ selected by false nearest neighbors.

Third, symmetry decides what can be reconstructed. A model-free cross-mapping test shows that the collar height is recoverable from the spring force ($\rho=0.99$) but not conversely ($\rho=-0.01$), because of the reflection $\theta\to-\theta$, and that the phase of the forcing restores identifiability ($\rho=0.99$). The mechanism is generic: a complete channel that is invariant under a symmetry of the dynamics, such as the $z$ variable of the Lorenz system under $(x,y)\to(-x,-y)$, cannot by itself identify a variable that the symmetry changes~\cite{LetellierAguirre2002}, and the same test detects the problem before any model is trained.

Fourth, recast as conditional density estimation with a neural spline flow, the FPRM gains most from its conditioning information and from its ability to represent several branches. With the spring force observed and $90\%$ of the collar-height and angular-velocity records missing, the physics-regularized flow lowers the continuous ranked probability score by a factor of 4.5 relative to the reimplemented original FPRM, by a factor of 3 when the latter is trained on every labeled sample, and by a factor of 1.8 relative to a Gaussian process with the same conditioning vector; it also recovers the fractal Poincar\'e section; when the target is ambiguous, the flow preserves both branches while regression averages them. At large missing fractions, physical residuals, holonomic or dynamical, help only when they are informative about the ambiguous degree of freedom and well conditioned; otherwise they can induce mode collapse.

Fifth, the complete study is available as an open laboratory, with a Python package, an interactive dashboard, and a browser-based simulation, in which every figure and number is regenerated from fixed seeds. Because its predictions can be checked against exact mechanics, the PIAR offers a controlled setting in which new reconstruction engines, physical priors, and sensor layouts can be tested before they are deployed on the incomplete records of real mechanical systems.
